\chapter{Guofeng: Leaving Breath Room in the Five Tones}
\label{ch:genre-guofeng-deep}

\classquote{大音希声，大象无形。}{《道德经》第四十一章}{The greatest sound seems faint; the greatest image has no outline. Guofeng's measure often lives in ``space'' and ``breath room'', not in piling up folk instruments.}

\section{What guofeng is: the flavour is more than a scale}
\label{sec:guofeng-what}

\term{guofeng}{guofeng} is not ``a pentatonic scale plus a flute''. Real guofeng blends the cadence of the Chinese language, the phrasing of opera, the playing techniques of folk instruments and contemporary arrangements into a sound you instantly recognise. It can be pure folk, or fused with electronic, rock and hip-hop, or a single operatic line inside a pop song. What makes it guofeng is breath, ornament and space---not how many Chinese symbols you use.

To hear guofeng, start with four handles:

\begin{longtable}{p{0.16\textwidth}p{0.32\textwidth}p{0.38\textwidth}}
\toprule
Handle & What it feels like & The writing question \\
\midrule
Breath room & Space between phrases, like a speaker pausing & Does this line leave room for an ornament to happen? \\
Ornament & Slides, grace notes and vibrato give a single note ``accent'' & Does the ornament grow out of the melody, or is it glued on every note? \\
Timbre & Flute, guzheng and erhu are distinctive without being loud & Are the folk instruments leads, or buried in the background? \\
Space & Less is more; empty space grabs more than fullness & Can this section lose one instrument and gain a beat of silence? \\
\bottomrule
\end{longtable}

The trap of guofeng is ``labelling'': a pentatonic scale and a few slides do not make a song guofeng. What moves people is often one line with real vocal inflection and one well-placed moment of restraint.

\section{From old scores to guofeng pop: how guofeng got here}
\label{sec:guofeng-history}

Guofeng is rooted in the gong-shang-jue-zhi-yu pentatonic system, poetic metre and operatic phrasing, and was reorganised in the twentieth century through recording, arrangement and fusion with Western orchestral and pop forms. The table follows one thread: when transmission changed, writing changed with it.

\begin{longtable}{p{0.16\textwidth}p{0.34\textwidth}p{0.42\textwidth}}
\toprule
Period & Key shift & Effect on arrangement \\
\midrule
\endfirsthead
\toprule
Period & Key shift & Effect on arrangement \\
\midrule
\endhead
Pre-Qin--Ming/Qing & Pentatonic scale, pitch pipes and poetry set to music; solo traditions for guqin, pipa, flute & Melodic contour and ``breath room'' come first; single-timbre solos with much space \\
Modern era & Records and radio; traditional tunes collected and arranged for ensembles & Folk ensemble sections form: bowed, plucked, blown and percussion \\
Mid--late 20th c. & National orchestras and symphonic arrangements & Folk instruments enter large ensembles; richer harmony and texture---and the risk of ``piling up'' \\
1990s--2000s & Pop industry and ``Chinese style'' songs & Pentatonic lines and operatic vocals enter pop; folk samples meet pop rhythm sections \\
2010s & Online music and independent artists & Guofeng/gufeng communities grow; freer mixes of electronic, hip-hop and folk \\
2020s--now & Short video, guochao and game/film scoring & Shorter sections, earlier hooks; folk timbres become cross-media cultural symbols \\
\bottomrule
\end{longtable}

Guofeng is a two-way street: traditional tunes are rearranged into contemporary works, and contemporary arranging in turn changes how folk instruments are played and recorded. When writing guofeng, ask first ``does this line sound like someone speaking?'', then ``which instrument says it?''

\section{Instrumentation: give every folk instrument a job}
\label{sec:guofeng-craft}

Guofeng lineups can be flexible, but ``who does what'' must be clear. Write functions into track names before stacking ten folk instruments.

\subsection{Typical lineup and four roles}
\label{sec:guofeng-instrumentation}

\begin{longtable}{p{0.16\textwidth}p{0.25\textwidth}p{0.25\textwidth}p{0.20\textwidth}}
\toprule
Layer & Common instruments/tracks & Main job & Common mistake \\
\midrule
\endfirsthead
\toprule
Layer & Common instruments/tracks & Main job & Common mistake \\
\midrule
\endhead
Melody & Dizi/xiao, erhu, lead/opera vocal & Carry the theme, give inflection, add ornament & Three at full volume at once, none of them a lead \\
Colour & Guzheng, pipa, yangqin & Glissando/arpeggios at section joints and phrase ends & Laid from start to finish, losing the ``point'' \\
Foundation & Pad, electronic bass, low strings & Hold harmony and low end & Fighting the drum/gong for the same low frequencies \\
Rhythm & Tanggu, wooden fish, bells, electronic drums & Mark sections, carry groove & Forcing a four-on-the-floor pattern that buries breath room \\
\bottomrule
\end{longtable}

The simplest start is five tracks: a lead (flute or voice), one plucked string (guzheng/pipa), one bowed string (erhu), one low layer (pad/low strings) and one percussion group. Get ``who speaks, who answers, who leaves space'' working before adding a second folk instrument or an electronic layer.

\subsection{Scale and melody: pentatonic is a skeleton, not a cage}
\label{sec:guofeng-scale}

The \term{pentatonic scale}{pentatonic-scale} (gong-shang-jue-zhi-yu) is guofeng's skeleton, not its limit. Qingjue and biangong colour tones bring the line closer to real folk singing; using only five notes can sound like an exercise. When writing a melody:

\begin{itemize}
\item Write a line that sounds like speech first, then decide which notes sit in the pentatonic frame;
\item Put ornaments (slides, grace notes, vibrato) on long notes and phrase ends, not every note;
\item Leave a beat of ``breath room'' at the phrase end for a guzheng glissando or a flute tail.
\end{itemize}

Pentatonic is not an excuse for simplicity. Writing five notes with a beginning, development, turn and return is harder---and more rewarding---than stacking ten.

\subsection{Groove and breath room: let the line breathe}
\label{sec:guofeng-groove}

Guofeng is often mid-tempo and unhurried; forcing a four-on-the-floor electronic beat buries the breath room. What works instead:

\begin{itemize}
\item \textbf{Percussion as section markers}: tanggu/big drum on section joints and strong beats, not every bar;
\item \textbf{Wooden fish/bells for fine grain}: mark breathing inside the bar instead of dense closed hats;
\item \textbf{Electronic drums half a beat back}: if you must use a modern rhythm section, offset kick/clap from the melody's accents to leave room for ornaments;
\item \textbf{Space at phrase ends}: let the melody finish and leave one beat with only pad or a flute tail---that beat is often where the guofeng feel lives.
\end{itemize}

\subsection{Harmony: transparency over complexity}
\label{sec:guofeng-harmony}

The keyword of guofeng harmony is transparency. Don't pile on sevenths and ninths; let the pentatonic melody land on clear chord tones. Practical moves:

\begin{itemize}
\item Use pentatonic tones as chord tones; colour tones only as passing notes;
\item Keep one owner of the low end: electronic bass and tanggu take turns; pad steps back;
\item At a turning point, use one unresolved cadence for ``more to come'', then return to the tonic in the reprise.
\end{itemize}

\subsection{Sound design: measure of ornament and space}
\label{sec:guofeng-sound}

Much of guofeng arranging is the placement of ornaments and space. Three principles:

\begin{itemize}
\item \textbf{Guzheng/pipa as points, not a bed}: glissandi and arpeggios at section joints and phrase ends, never from start to finish.
\item \textbf{Flute forward, keep the breath}: the breath and slight airiness of the lead are the ``human'' part---don't over-clean them.
\item \textbf{Erhu as an answer}: in the A section let erhu answer the flute/voice theme, like a second person replying.
\end{itemize}

Arrange ``ornament--space--answer'' well and you need fewer folk instruments to sound like guofeng.

\section{Four pieces, arranged}
\label{sec:guofeng-cases}

The four pieces below are arrangement analyses only---no lyrics, no full melodies. Work information is sourced from public references and official pages; links are in Further Reading.

\subsection{``Jasmine Flower'': the breath room of a folk song}
\label{sec:guofeng-case-jasmine}

``Jasmine Flower'' (Molihua) is a folk song from the Jiangsu area, widely known after being collected and rearranged, and often used in film and international performances. Its power is not complexity but the breathing of its phrases.

\begin{itemize}
\item \textbf{Melody}: a pentatonic skeleton with natural drag at phrase ends; every phrase leaves ``breath room''.
\item \textbf{Arrangement}: versions range from voice or a single folk instrument to full orchestra, all keeping the melody clear and the background non-competing.
\item \textbf{Arrangement logic}: the simpler the lineup, the more the melody is tested---added instruments must step back for the lead.
\end{itemize}

Takeaway: when rearranging a folk song, preserve its phrasing and space before thickening the texture.

\subsection{``Spring River Flower Moon Night'': a pipa tune unfolds}
\label{sec:guofeng-case-spring-river}

``Spring River Flower Moon Night'' was adapted from the pipa piece ``Sunset Flute and Drum'' and arranged for ensemble by Shanghai's Datong Music Society in the 1920s. It shows how ``beginning--development--turn--return'' unfolds through arrangement.

\begin{itemize}
\item \textbf{Theme}: the core melody is played by different instruments across sections---one painting seen from different angles.
\item \textbf{Texture}: plucked, blown and bowed groups take turns on the theme; sections join through changes of tempo and density.
\item \textbf{Arrangement logic}: not a pile of theme-and-variations, but timbre and thickness following the narrative.
\end{itemize}

Takeaway: for a long guofeng section, let one theme pass through three instruments instead of writing three new melodies.

\subsection{``Blue and White Porcelain'': the grammar of Chinese-style pop}
\label{sec:guofeng-case-blue-and-white}

``Blue and White Porcelain'' (Qinghua Ci), lyrics by Vincent Fang, arrangement by Baby Chung and music by Jay Chou, was released on 2 November 2007 on the album \emph{On the Run!}. It is a representative ``Chinese style'' pop song in the mainstream.

\begin{itemize}
\item \textbf{Melody and words}: pentatonic motion follows Mandarin tones and imagery; short phrases repeat with subtle changes.
\item \textbf{Lineup}: a pop rhythm section supports; pipa/guzheng colour decorates without covering the lead vocal.
\item \textbf{Arrangement logic}: guofeng is ``grammar'' rather than ``scenery''---it serves the lyrics' images and emotions.
\end{itemize}

Takeaway: treat folk instruments as a way of speaking, not decorations pinned on; let them tell the story with the lyrics.

\subsection{``Red Costume'': operatic vocal meets contemporary production}
\label{sec:guofeng-case-chiling}

``Red Costume'' (Chiling), sung by HITA and released in 2018, became widely popular among young listeners for mixing operatic vocal style with modern arrangement. It shows how traditional opera elements enter contemporary production.

\begin{itemize}
\item \textbf{Vocal}: the operatic section uses sung-and-spoken inflection and ornament, contrasting sharply with the pop verses.
\item \textbf{Lineup}: electronic and folk timbres coexist; sections advance through changes in vocal form and density.
\item \textbf{Arrangement logic}: the opera line carries dramatic conflict---the arrangement makes space for it rather than filling it.
\end{itemize}

Takeaway: before adding an operatic line, decide what emotion it carries; giving it space beats giving it more tracks.

\section{Yeban sketch: eight bars of guofeng}
\label{sec:guofeng-yeban}

An eight-bar sketch you can rebuild in Yeban, using ``flute/voice + plucked + bowed + low layer + percussion'' to run guofeng's breath room.

\begin{longtable}{p{0.10\textwidth}p{0.22\textwidth}p{0.58\textwidth}}
\toprule
Bars & Section & Arrangement moves \\
\midrule
\endfirsthead
\toprule
Bars & Section & Arrangement moves \\
\midrule
\endhead
1--4 & A (narration) & Dizi/lead: pentatonic theme, a slide on the phrase end and one beat of breath room; guzheng: a short glissando at the phrase end, not a bed; erhu: answers one phrase a fifth below; pad: mid-low sustained tones leaving the breath room clear; percussion: light wooden fish, tanggu once in bar 1 \\
5--8 & B (development) & Melody: theme an octave higher with one vibrato; guzheng: denser arpeggios but still leaving phrase ends empty; electronic bass: enters, alternating the low end with tanggu; erhu: long bow underneath, not competing; bar 8 drops all percussion, leaving only flute tail and pad \\
\bottomrule
\end{longtable}

Build order:

\begin{sloppypar}
\begin{enumerate}
\item Create a mid-tempo (about 84 BPM), 4/4 project and build tracks as above: Vocal/Dizi (MIDI or audio), Guzheng (MIDI), Erhu (MIDI), Pad (MIDI), Percussion (MIDI), Master bus.
\item Write the flute/voice pentatonic theme first: two four-bar phrases with a slide and a beat left at each end; solo-listen to confirm it ``sounds like speaking''.
\item Write guzheng glissandi/arpeggios only at phrase ends and section joints, never from start to finish.
\item Write erhu's answer in the A section with a clear landing; switch to long bow underneath in B.
\item Write pad and low layer: mid-low pad in A; in B electronic bass and tanggu alternate the low end---check levels (F8) so they don't fight.
\item Write percussion: wooden fish/bells for breathing, tanggu on strong beats and section points; drop it all in bar 8.
\item Open the piano roll (F7) to place ornaments; open the device stack (F3) so no two folk instruments fight for one register.
\item Save and export an eight-bar listen: if you hear ``narration--development--space--return'' and no instrument plays from start to finish, it works.
\end{enumerate}
\end{sloppypar}

\section{Exercises: three routes, different trade-offs}
\label{sec:guofeng-practice}

\begin{exercise}
\textbf{One pentatonic theme, three ornament densities.} Write an 8-bar pentatonic line in three versions: A slides only at phrase ends; B a grace note every two notes; C vibrato throughout. Compare inflection, clarity and listenability. A sounds most like speech but may feel thin; C has the most ``flavour'' but cloys; B sits between. Ask blind listeners which sounds ``sung'' rather than ``played as a scale''.
\end{exercise}

\begin{exercise}
\textbf{Where does the breath room go?} Make three versions of the same line: A one beat at phrase ends; B half a beat mid-phrase; C no space. Play each and note which makes you want the next phrase. Then remove the pad in the empty space, leaving only the folk tail---hear whether the ``emptiness'' grabs more.
\end{exercise}

\begin{exercise}
\textbf{One folk instrument, three roles.} Use only guzheng and make it: A state the theme; B glissando at section joints; C a low bed throughout. Compare which is least intrusive and which feels most like ``the point'' of guofeng. Don't add a second folk instrument---use division of labour only.
\end{exercise}

\section{Further reading}
\label{sec:guofeng-reading}

\begin{sloppypar}
\begin{itemize}
\item \textbf{Wikipedia, ``Mo Li Hua'':} \url{https://en.wikipedia.org/wiki/Mo_Li_Hua}. The ``Jasmine Flower'' folk song---its circulation, regional variants and adaptations, the source of guofeng melody.
\item \textbf{Wikipedia, ``The Moon over the River on a Spring Night'':} \url{https://en.wikipedia.org/wiki/The_Moon_over_the_River_on_a_Spring_Night}. Zhang Ruoxu's Tang poem, from which the piece takes its name---the literary background to its imagery.
\item \textbf{Wikipedia, ``Blue and White Porcelain'':} \url{https://en.wikipedia.org/wiki/Blue_and_White_Porcelain}. The Jay Chou/Vincent Fang entry, with writer, arranger and release details.
\item \textbf{Wikipedia, ``Zhongguo feng (music)'':} \url{https://en.wikipedia.org/wiki/Zhongguo_feng_(music)}. An overview of Chinese-style/guofeng pop, its contemporary development and young-listener context.
\end{itemize}
\end{sloppypar}

The links above were checked against entry names and release data on 27 November 2026 (work/writer/date are in each entry's text); Wikipedia URLs are stable---if your region blocks them, search the same title in a local language edition. For melody and timbre, listen on the artists' official channels by title.

\section{Takeaways}
\label{sec:guofeng-takeaways}

\begin{itemize}
\item Guofeng's flavour lives in breath, ornament and space---not ``pentatonic scale plus a flute''.
\item The scale is a skeleton, not a cage: qingjue/biangong for colour, ornaments on long notes and phrase ends.
\item Write functions first: melody, colour, foundation and rhythm each have a job; more folk instruments are not better.
\item Breath room is guofeng's respiration: space at phrase ends, percussion as section markers, electronic drums yielding to ornaments.
\item Seek transparent harmony over complexity; one owner of the low end, tanggu and electronic bass taking turns.
\item Guzheng/pipa as points not a bed, flute forward with its breath, erhu as an answer.
\item Build the Yeban sketch in the order ``theme--ornament--answer--low layer--percussion--space/return''.
\end{itemize}

\begin{helpmapping}
\begin{longtable}{p{0.45\textwidth}p{0.45\textwidth}}
\toprule
Section & Help anchor \\
\midrule
What guofeng is & \texttt{arrangement-genre-guofeng-what} \\
From old scores to guofeng pop & \texttt{arrangement-genre-guofeng-history} \\
Instrumentation: every folk instrument has a job & \texttt{arrangement-genre-guofeng-craft} \\
Four pieces, arranged & \texttt{arrangement-genre-guofeng-cases} \\
Yeban sketch & \texttt{arrangement-genre-guofeng-yeban} \\
Exercises & \texttt{arrangement-genre-guofeng-practice} \\
Further reading & \texttt{arrangement-genre-guofeng-reading} \\
\bottomrule
\end{longtable}
\end{helpmapping}

\begin{figure}[htbp]\centering
\begin{tikzpicture}[font=\scriptsize,x=1cm,y=1cm,>=Stealth]
\yebangenreddraw{Guofeng: bring traditional timbres into a modern arrangement}{60--200 Hz drums and bass}{200--2 kHz dizi / pipa / voice}{2--10 kHz plucked strings and air}
\end{tikzpicture}
\caption{Guofeng frequency split, drum grid and section-energy curve}
\end{figure}
